% Source printed pp. 136-141, physical PDF pp. 144-149.
\section{维数, 切空间与剩余 gerbe}\label{src:4.5-section}
\begin{epigraph}
点是没有部分的东西.
\epigraphauthor{Euclid}
\end{epigraph}
将维数和切空间的概念推广到代数叠之后, 我们证明每个有限型点都有一个剩余 gerbe, 用来替代概形的点的剩余域这一概念.
\subsection{维数}\label{src:4.5.1-subsection}
回忆, 概形 $X$ 的维数 $\dim X$ 是其底层拓扑空间的 Krull 维数, 而点 $x\in X$ 处的维数 $\dim_x X$ 是包含 $x$ 的开子集的维数的最小值. 当 $X$ 在域上有限型且 $x\in X$ 是闭点时, 它等于 $\dim\mathcal O_{X,x}$. 现在将这些定义推广到代数空间和叠.
\begin{definition}[维数]\label{src:4.5.1}
\begin{enumerate}[(1)]
\item 设 $X$ 是 Noether 代数空间, $x\in|X|$. 定义 $X$ 在 $x$ 处的\newterm{维数}{dimension} 为
\[
\dim_x X=\dim_u U\in\mathbb Z_{\geq0}\cup\infty,
\]
其中 $U\to X$ 是平展呈示, $u\in U$ 是 $x$ 的一个原像.
\item 设 $\mathcal X$ 是 Noether 代数叠, $x\in|\mathcal X|$. 设 $U\to\mathcal X$ 是光滑呈示, $s,t:R\rightrightarrows U$ 是所诱导的光滑群胚, $u\in U$ 是 $x\in|\mathcal X|$ 的一个原像. 定义 $\mathcal X$ 在 $x$ 处的维数为
\[
\dim_x\mathcal X=\dim_u U-\dim_{e(u)}R_u\in\mathbb Z\cup\infty,
\]
其中 $R_u$ 是 $s:R\to U$ 在 $u$ 上的纤维, $e:U\to R$ 表示群胚中的单位态射.
\item 若 $\mathcal X$ 是 Noether 代数空间或叠, 则定义 $\mathcal X$ 的维数为
\[
\dim\mathcal X=\sup_{x\in|\mathcal X|}\dim_x\mathcal X\in\mathbb Z\cup\infty.
\]
\end{enumerate}
\end{definition}
\begin{proposition}\label{src:4.5.2}
$\dim_x\mathcal X$ 的定义与呈示 $U\to\mathcal X$ 及 $x$ 的原像 $u$ 的选择无关.
\end{proposition}
\begin{proof}
由于平展态射的相对维数为 $0$, 代数空间在一点处的维数是良定义的. 若 $U\to\mathcal X$ 是光滑呈示 (其中 $U$ 是概形), $u\in U$ 是 $x$ 的一个原像, 剩余域为 $\kappa(u)$, 则纤维 $R_u$ 可等同于下列纤维积:
\[
\begin{tikzcd}
R_u\arrow[r]\arrow[d]&R\arrow[r,"t"]\arrow[d,"s"]&U\arrow[d]\\
\Spec\kappa(u)\arrow[r]\arrow[ur,phantom,"\square" description]&U\arrow[r]\arrow[ur,phantom,"\square" description]&\mathcal X,
\end{tikzcd}
\]
它是 $\kappa(u)$ 上的光滑代数空间. 设 $U'\to\mathcal X$ 是另一呈示, 诱导光滑群胚 $R'\rightrightarrows U'$, 且 $u'\in U'$ 是 $x$ 的一个原像. 由于可构造同时支配 $U$ 和 $U'$ 的光滑呈示 $U\times_{\mathcal X}U'$, 可以假设存在将 $u'$ 送到 $u$ 的概形的光滑态射 $U'\to U$. 由光滑态射的维数公式,
\begin{equation}\label{src:4.5.3}
\dim_{u'}U'=\dim_u U+\dim_{u'}U'_u;
\end{equation}
见命题~\bookref{A.3.10}. 为了将 $\dim_{u'}U'_u$ 与 $\dim_{e(u)}R_u$ 和 $\dim_{e'(u')}R'_{u'}$ 联系起来, 构造下列 Cartesian 图:
\[
\begin{tikzcd}[column sep=small]
R'_{u'}\arrow[r]\arrow[d]&U''_u\arrow[r]\arrow[d]&U''\arrow[r]\arrow[d]&U'\arrow[d]\\
R_u\times_{\kappa(u)}\kappa(u')\arrow[r]\arrow[d]\arrow[ur,phantom,"\square" description]&R_u\arrow[r]\arrow[d]\arrow[ur,phantom,"\square" description]&R\arrow[r,"t"]\arrow[d,"s"]\arrow[ur,phantom,"\square" description]&U\arrow[d]\\
\Spec\kappa(u')\arrow[r]\arrow[ur,phantom,"\square" description]&\Spec\kappa(u)\arrow[r]\arrow[ur,phantom,"\square" description]&U\arrow[r]\arrow[ur,phantom,"\square" description]&\mathcal X.
\end{tikzcd}
\]
再次对 $U''_u\to R_u$ 应用光滑态射的维数公式, 但这一次使用其对代数空间的推广 (习题~\bookref{4.5.5}(a)). 利用 $e(u)\in|R_u|$ 上的纤维可等同于 $U'_u$, 并取 $u''\in U''_u$ 为 $e(u)\in|R_u|$ 和 $u'\in U'$ 的一个原像, 得到
\begin{equation}\label{src:4.5.4}
\dim_{u''}U''_u=\dim_{e(u)}R_u+\dim_{u'}U'_u.
\end{equation}
由于维数在域扩张下不变 (习题~\bookref{4.5.5}(b)), 有 $\dim_{e'(u')}R'_{u'}=\dim_{u''}U''_u$. 结合~\bookeqref{4.5.3} 和~\bookeqref{4.5.4} 即得所需结论.
\end{proof}
\begin{exercise}[细节]\label{src:4.5.5}
\begin{enumerate}[(a)]
\item 证明, 若 $X\to Y$ 是 Noether 代数空间的光滑态射, $x\in|X|$ 的像为 $y\in|Y|$, 则
\[
\dim_x X=\dim_y Y+\dim_x X_y.
\]
这将命题~\bookref{A.3.10} 推广到代数空间.
\item 设 $\mathcal X$ 是域 $k$ 上的 Noether 代数叠, $k\to k'$ 是域扩张, $x'\in|\mathcal X_{k'}|$ 的像为 $x\in|\mathcal X|$. 证明 $\dim_{x'}\mathcal X_{k'}=\dim_x\mathcal X$.
\end{enumerate}
\end{exercise}
\begin{example}\label{src:4.5.6}
若 $G$ 是域 $k$ 上的光滑仿射代数群, 作用于 $k$ 上有限型的纯维概形 $U$, 则
\[
\dim[U/G]=\dim U-\dim G.
\]
特别地, 分类叠 $BG$ 的维数为
\[
\dim BG=-\dim G,
\]
它可以为负! 类似地, $\dim[\mathbb A^1_k/\mathbb G_m]=0$.
\end{example}
\subsection{切空间}\label{src:4.5.2-subsection}
\newterm{对偶数}{dual numbers} 是域 $k$ 上的环 $k[\epsilon]:=k[\epsilon]/\epsilon^2$.
\begin{definition}\label{src:4.5.7}
对于代数叠 $\mathcal X$ 和点 $x:\Spec k\to\mathcal X$, 定义 \newterm{Zariski 切空间}{Zariski tangent space} (或简称 $\mathcal X$ 在 $x$ 处的切空间) 为集合
\[
T_{\mathcal X,x}:=\left\{
\text{$2$-交换图}\quad
\begin{tikzcd}[column sep=2.8em,row sep=3em]
\Spec k\arrow[d,hook]\arrow[dr,"x",""{name=x}]&{}\\
\Spec k[\epsilon]\arrow[r,"\tau"',""{name=tau}]&\mathcal X
\arrow[from=x,to=tau,Rightarrow,"\alpha"']
\end{tikzcd}
\right\}\big/\sim.
\]
换言之, $T_{\mathcal X,x}$ 是二元组 $(\tau,\alpha)$ 的集合, 其中 $\tau:\Spec k[\epsilon]\to\mathcal X$, $\alpha:x\xrightarrow{\sim}\tau|_k$. 两个二元组等价 $(\tau,\alpha)\sim(\tau',\alpha')$, 当且仅当在 $\mathcal X(k[\epsilon])$ 中存在与 $\alpha$ 和 $\alpha'$ 相容的同构 $\beta:\tau\xrightarrow{\sim}\tau'$, 即 $\alpha'=\beta|_{\Spec k}\circ\alpha$.
\end{definition}
\begin{example}\label{src:4.5.8}
考虑定义在 $k$ 上, 亏格 $g\geq2$ 的光滑连通射影曲线 $[C]\in\mathcal M_g(k)$. 形变理论~\bookeqref{C.1.11} 给出 $T_{\mathcal M_g,[C]}=H^1(C,T_C)$. 由于 $\deg T_C<0$, 有 $H^0(C,T_C)=0$, 而 Riemann--Roch 定理给出
\[
\begin{aligned}
\dim T_{\mathcal M_g,[C]}&=\dim H^1(C,T_C)=-\chi(T_C)\\
&=-\bigl(\deg T_C+(1-g)\bigr)=3g-3.
\end{aligned}
\]
\end{example}
\begin{example}\label{src:4.5.9}
设 $C$ 是 $k$ 上的光滑连通射影曲线, $E\in\cBun_{r,d}(C)(k)$ 是 $C$ 上秩为 $r$, 次数为 $d$ 的向量丛. 形变理论~\bookeqref{C.1.18} 给出 $T_{\cBun_{r,d}(C),[E]}=\operatorname{Ext}^1_{\mathcal O_C}(E,E)=H^1(C,E\otimes E^\vee)$. 由 Riemann--Roch 定理~\bookeqref{9.1.2}, 有 $\chi(E\otimes E^\vee)=r^2(1-g)$. 由于 $\dim\Aut(E)=\dim_k\operatorname{Hom}_{\mathcal O_C}(E,E)=H^0(C,E\otimes E^\vee)$, 算得
\[
\dim T_{\cBun_{r,d}(C),[E]}=\dim\operatorname{Ext}^1_{\mathcal O_C}(F,F)=\dim\Aut(F)+r^2(g-1).
\]
\end{example}
\begin{proposition}\label{src:4.5.10}
若 $\mathcal X$ 是具有仿射对角线的代数叠, $x\in\mathcal X(k)$, 则 $T_{\mathcal X,x}$ 自然是一个 $k$-向量空间.
\end{proposition}
\begin{proof}
$c\in k$ 对 $(\tau,\alpha)\in T_{\mathcal X,x}$ 的数乘由复合 $\Spec k[\epsilon]\to\Spec k[\epsilon]\xrightarrow{\tau}\mathcal X$ 定义, 其中第一个映射由 $\epsilon\mapsto c\epsilon$ 给出, 并使用同一个 $2$-同构 $\alpha$. 为了定义加法, 我们证明存在范畴等价
\[
\mathcal X(k[\epsilon_1]\times_k k[\epsilon_2])\longrightarrow\mathcal X(k[\epsilon_1])\times_{\mathcal X(k)}\mathcal X(k[\epsilon_2]),
\]
或者换言之, 下图在具有仿射对角线的代数叠中是推出图:
\[
\begin{tikzcd}
\Spec k\arrow[r,hook]\arrow[d,hook]\arrow[dr,phantom,"\ulcorner" description]&\Spec k[\epsilon_1]\arrow[d,hook]\\
\Spec k[\epsilon_2]\arrow[r,hook]&\Spec(k[\epsilon_1]\times_k k[\epsilon_2]).
\end{tikzcd}
\]
(这与 Schlessinger 判据~\bookeqref{C.4.6} 和 Artin 代数性公理~\bookeqref{C.7.4} 中的齐性条件有关; 概形推出的讨论另见第~\bookref{B.4}~节.) 建立这一结论后, 通过复合 $\Spec k[\epsilon]\to\Spec(k[\epsilon_1]\times_k k[\epsilon_2])\to\mathcal X$ 定义 $(\tau_1,\alpha_1)$ 与 $(\tau_2,\alpha_2)$ 的加法, 其中第一个映射由将 $(\epsilon_1,0)$ 和 $(0,\epsilon_2)$ 都送到 $\epsilon$ 来定义.

为了建立这一等价, 选取从仿射概形 $U$ 出发的光滑态射 $(U,u)\to(\mathcal X,x)$. 由于 $\mathcal X$ 具有仿射对角线, $U\to\mathcal X$ 是仿射态射. 令 $\Spec A_0=\Spec k\times_{\mathcal X}U$, $\Spec A_1=\Spec k[\epsilon_1]\times_{\mathcal X}U$, $\Spec A_2=\Spec k[\epsilon_2]\times_{\mathcal X}U$. 由于 $\Spec(A_1\times_A A_2)$ 是仿射概形范畴中 $\Spec A_0\hookrightarrow\Spec A_1$ 和 $\Spec A_0\hookrightarrow\Spec A_2$ 的推出, 存在唯一的态射 $\Spec(A_1\times_A A_2)\to\Spec(k[\epsilon_1]\times_k k[\epsilon_2])$ 和 $\Spec(A_1\times_A A_2)\to U$, 补全下图:
\[
\begin{tikzcd}[column sep=small,row sep=small]
{}&\Spec A_0\arrow[rrr]\arrow[dd]\arrow[dl]&{}&{}&\Spec A_2\arrow[dd]\arrow[dl]&{}\\
\Spec k\arrow[rrr]\arrow[dd]&{}&{}&\Spec k[\epsilon_1]\arrow[dd]\arrow[dddr,bend left=15,"\tau_1" near start]&{}&{}\\
{}&\Spec A_1\arrow[rrr]\arrow[dl]&{}&{}&\Spec(A_1\times_A A_2)\arrow[dl]\arrow[dr]&{}\\
\Spec k[\epsilon_2]\arrow[rrr]\arrow[drrrr,bend right=15,"\tau_2"']&{}&{}&\Spec(k[\epsilon_1]\times_k k[\epsilon_2])\arrow[dr,dashed]&{}&U\arrow[dl]\\
{}&{}&{}&{}&\mathcal X&{}
\end{tikzcd}
\]
由 Artin 环上的平坦性判据~\bookeqref{A.2.3}, 可知映射 $\Spec(A_1\times_A A_2)\to\Spec(k[\epsilon_1]\times_k k[\epsilon_2])$ 忠实平坦. 对 $U\times_{\mathcal X}U$ 重复这一论证, 可知 $\Spec(A_1\times_A A_2)\to U$ 唯一下降为映射 $U\to\mathcal X$. 这些细节, 以及数乘和加法赋予 $T_{\mathcal X,x}$ 向量空间结构的验证, 留给读者.
\end{proof}
\begin{exercise}[较难]\label{src:4.5.11}
\begin{enumerate}[(a)]
\item 证明 $T_{\mathcal X,x}$ 自然是 $G_x$ 的一个表示, 在集合层面由下式给出: 对 $g\in G_x$ 和 $(\tau,\alpha)\in T_{\mathcal X,x}$, 有 $g\cdot(\tau,\alpha)=(\tau,g\circ\alpha)$.
\item 证明命题~\bookref{4.5.10} 中对角线仿射的条件是多余的.
\end{enumerate}
\end{exercise}
\subsection{剩余 gerbe}\label{src:4.5.3-subsection}
概形的每个点 $x\in X$ 都有一个剩余域 $\kappa(x)$, 以及一个像为 $x$ 的单态射 $\Spec\kappa(x)\to X$. 剩余 gerbe 将为代数叠提供类似的性质. 注意, 非平凡稳定子使域值点无法成为单态射, 例如对于非平凡代数群, $\Spec k\to BG$ 不是单态射.
\begin{definition}\label{src:4.5.12}
设 $\mathcal X$ 是代数叠, $x\in|\mathcal X|$. 若存在约化的局部 Noether 代数叠 $\mathcal G_x$ 及单态射 $\mathcal G_x\hookrightarrow\mathcal X$, 使 $|\mathcal G_x|$ 是映到 $x$ 的单点, 则称 $x$ 处的\newterm{剩余 gerbe}{residual gerbe} 存在. 代数叠 $\mathcal G_x$ 称为 $x$ 处的剩余 gerbe.
\end{definition}
我们将在定义~\bookref{7.4.6} 中定义 gerbe, 并证明 $\mathcal G_x$ 是域 $\kappa(x)$ 上的 gerbe, 这个域称为\newterm{剩余域}{residue field} (命题~\bookref{7.4.39}). 虽然定义中看不出来, 但不难证明剩余 gerbe 实际上是唯一的 (见命题~\bookref{4.5.16}). 这些事实说明了我们称 $\mathcal G_x$ 为剩余 gerbe 的理由. 我们将在第~\bookref{7.4.5-subsection}~小节中更一般地讨论剩余 gerbe 的存在性; 目前则只验证有限型点处的剩余 gerbe 的存在性.
\begin{definition}\label{src:4.5.13}
代数叠中的点 $x\in|\mathcal X|$ 称为\newterm{有限型点}{finite type point}, 若存在 $x$ 的一个代表 $\Spec k\to\mathcal X$, 它是局部有限型的.\footnote{若 $\mathcal X$ 具有拟紧对角线, 例如 $\mathcal X$ 拟分离 (定义~\bookref{4.3.11}), 则每个域值点 $\Spec k\to\mathcal X$ 自动拟紧, 因而 $\Spec k\to\mathcal X$ 局部有限型等价于它有限型.}
\end{definition}
\begin{remark}\label{src:4.5.14}
若 $X$ 是 Noether 概形, 则点 $x\in X$ 有限型当且仅当 $x\in X$ 局部闭. 更一般地, 从域出发的态射 $\Spec k\to X$, 像为 $x$, 有限型当且仅当像 $x\in X$ 局部闭且 $\kappa(x)/k$ 是有限扩张. 事实上, 为了证明非平凡的方向 $(\Rightarrow)$, 将 $X$ 替换为 $\overline{\{x\}}$. 由于 $\Spec k\to X$ 有限型且像稠密, 泛平坦性~\bookeqref{A.2.13} 表明 $\Spec k\to X$ 是 fppf 的, 因而它的像是开的.

概形中有限型但不闭的点的一个例子是 DVR 的泛点. 另一方面, 若 $X$ 是域 $k$ 上的有限型概形, 则每个有限型点都是闭点. 对 $k$ 上的有限型代数叠, 类似的事实并不成立, 例如 $1:\Spec k\to[\mathbb A^1/\mathbb G_m]$ 是一个开的有限型点.
\end{remark}
\begin{exercise}\label{src:4.5.15}
设 $\mathcal X$ 是代数叠.
\begin{enumerate}[(a)]
\item 证明点 $x\in|\mathcal X|$ 有限型当且仅当存在概形 $U$, 闭点 $u\in U$ 和光滑态射 $(U,u)\to(\mathcal X,x)$.
\item 证明任意代数叠 (相应地, 拟紧代数叠) 都有一个有限型点 (相应地, 闭点).
\end{enumerate}
\end{exercise}
\begin{proposition}[剩余 gerbe 的存在性]\label{src:4.5.16}
若 $\mathcal X$ 是 Noether 代数叠, $x\in|\mathcal X|$ 是有限型点, 则 $x$ 处存在唯一的剩余 gerbe $\mathcal G_x$. 此外, $\mathcal G_x$ 满足:
\begin{enumerate}[(1)]
\item 代数叠 $\mathcal G_x$ 正则, 态射 $\mathcal G_x\hookrightarrow\mathcal X$ 是局部闭浸入.
\item 若进一步假设 $\mathcal X$ 在域 $k$ 上有限型, 且 $x\in\mathcal X(k)$ 的稳定子 $G_x$ 光滑且仿射, 则 $\mathcal G_x\cong BG_x$.
\item 若进一步假设 $\mathcal X$ 是 Noether 代数空间, 则对于某个称为 $x$ 的剩余域的域 $\kappa(x)$, 有 $\mathcal G_x\cong\Spec\kappa(x)$.
\end{enumerate}
\end{proposition}
\begin{proof}
将 $\mathcal X$ 替换为 $\overline{\{x\}}$ 后, 可以假设 $\mathcal X$ 约化且 $x\in|\mathcal X|$ 稠密. 设 $\Spec k\to\mathcal X$ 是 $x$ 的一个有限型代表. 由泛平坦性~\bookeqref{4.3.32}, $\Spec k\to\mathcal X$ 平坦. 因而 $\Spec k\to\mathcal X$ 的像 $\{x\}\subseteq|\mathcal X|$ 是开的 (习题~\bookref{4.3.27}). 相应的开子叠 $\mathcal G_x\subseteq\mathcal X$ 满足剩余 gerbe 的各项性质. 由于 $\Spec k\to\mathcal G_x$ 是 fppf 的, 而正则性沿 fppf 态射下降 (命题~\bookref{3.1.21}), $\mathcal G_x$ 正则.

对于唯一性, 假设 $\mathcal G$ 和 $\mathcal G'$ 是约化的局部 Noether 代数叠, 带有单态射 $\mathcal G\hookrightarrow\mathcal X$ 和 $\mathcal G'\hookrightarrow\mathcal X$, 且 $|\mathcal G|$ 和 $|\mathcal G'|$ 都是映到 $x$ 的单点集. 那么 $\mathcal G'':=\mathcal G\times_{\mathcal X}\mathcal G'$ 是非空代数叠, 带有单态射 $\mathcal G''\to\mathcal G$ 和 $\mathcal G''\to\mathcal G'$. 取从域出发的有限型态射 $\Spec k\to\mathcal G$, 它由习题~\bookref{4.5.15} 保证存在, 且由泛平坦性~\bookeqref{4.3.32} 是 fppf 的. 基变换 $\mathcal G''\times_{\mathcal G}\Spec k$ 是非空的 Noether 代数空间, 带有到 $\Spec k$ 的单态射. 由习题~\bookref{4.5.17}, 这蕴含 $\mathcal G''\times_{\mathcal G}\Spec k\cong\Spec k$. 由 fppf 下降, $\mathcal G''\to\mathcal G$ 也是同构. 类似地, $\mathcal G''\to\mathcal G'$ 是同构.

对于 (2), 存在预叠的单态射 $BG_x^{\mathrm{pre}}\to\mathcal X$: 对于 $k$-概形 $T$, $BG_x^{\mathrm{pre}}$ 在 $T$ 上有唯一的对象, 将此对象映到复合 $T\to\Spec k\xrightarrow{x}\mathcal X$. 类似地, 位于 $T'\to T$ 上的态射对应于一个映射 $T'\to G_x$, 将其映到 $\mathcal X$ 中相应的态射. 叠化后, 这诱导单态射 $BG_x\to\mathcal X$, 而且可以看到 $BG_x$ 满足剩余 gerbe 的各项性质. 最后, (3) 由习题~\bookref{4.5.17} 推出.
\end{proof}
\begin{exercise}[细节]\label{src:4.5.17}
设 $X$ 是约化的 Noether 代数空间, 且 $|X|$ 是一个点. 证明对于某个域 $k$, 有 $X\cong\Spec k$. (这也将由后面的事实推出, 例如定理~\bookref{5.5.1}.)
\end{exercise}
\begin{exercise}[技术性]\label{src:4.5.18}
设 $\mathcal X$ 是代数叠 (可以不是 Noether 的), $x\in\mathcal X$ 是有限型点, 其稳定子非分歧 (即任意代表的稳定子群概形都非分歧). 证明剩余 gerbe 存在且唯一. 另见 \cite[Tag 06G3]{SP}.
\end{exercise}
下面的推论说明, 光滑群胚中有限型点的轨道 $O(u)\subseteq U$ 是局部闭的.
\begin{corollary}\label{src:4.5.19}
设 $x\in|\mathcal X|$ 是 Noether 代数叠 $\mathcal X$ 的有限型点. 若 $(U,u)\to(\mathcal X,x)$ 是从概形 $U$ 出发的光滑态射, 且 $u\in U$ 是有限型点, 则存在 Cartesian 图
\begin{equation}\label{src:4.5.20}
\begin{tikzcd}
O(u)\arrow[r,hook]\arrow[d]&U\arrow[d]\\
\mathcal G_x\arrow[r,hook]\arrow[ur,phantom,"\square" description]&\mathcal X
\end{tikzcd}
\end{equation}
其中 $O(u)$ 在集合层面可等同于所诱导群胚 $s,t:R:=U\times_{\mathcal X}U\rightrightarrows U$ 的轨道 $s(t^{-1}(u))$.
\end{corollary}
\begin{samepage}
\begin{remark}\label{src:4.5.21}
若 $\mathcal X=[U/G]$ 是域 $k$ 上的光滑仿射代数群作用于 Noether $k$-概形 $U$ 所得的商叠, 且 $u\in U(k)$, 则存在 Cartesian 图
\[
\begin{tikzcd}
Gu\arrow[r,hook]\arrow[d]&U\arrow[d]\\
BG_x\arrow[r,hook]\arrow[ur,phantom,"\square" description]&{[U/G]},
\end{tikzcd}
\]
于是得到熟知的事实: 轨道 $Gu\hookrightarrow U$ 是局部闭的 (命题~\bookref{B.1.20}(5)).
\end{remark}
\end{samepage}
\begin{exercise}\label{src:4.5.22}
设 $\mathcal X$ 是 Noether 代数叠, $x\in|\mathcal X|$ 是稳定子光滑的有限型点. 设 $\overline x:\Spec k\to\mathcal X$ 是 $x$ 的一个代表. 证明 $\dim\mathcal G_x=-\dim G_{\overline x}$.
\end{exercise}
